\scanpage{23}
\mspage{11}
4. \uline{Relations avec le cube orienté}

Si $I$ est un ens. à $4$ éléments, $J\simeq\mathcal{P}_{2,2}(I)\simeq V_I^*$ l'ens à trois éléments associé, on considère l'application de degré $2$
\diagram{\begin{tikzcd}
(*)\quad \operatorname{Circ}(I) \arrow[r] & J\\
\substack{\text{ens des permutations}\\\text{circulaires sur }I} &\\
u \arrow[r,mapsto] & u^2
\end{tikzcd}}
[NB On a aussi l'application de degré $2$
\diagram{\begin{tikzcd}
\mathcal{P}_2(I) \arrow[r] & J\\
A \arrow[r,mapsto] & \{A,I\setminus A\}
\end{tikzcd}}
$\simeq$ ens. des ``\unsure{dicompositions|décompositions} de $I$'' (ou ``réflexions'')
\[
\text{ou }\forall\tau\ \text{réflexion},\
\exists!\tau'\ \text{réflexion telle que }\tau\tau'\in V^*.
\]
ce sont vraiment des revêtements de nature distincte, les espaces homogènes \inserted{sous $G$} $\operatorname{Circ}(I)$ et $\mathcal{P}_2(I)$ n'étant pas \unsure{$\tilde m$} isomorphes (le stabilisateur d'un élément de $\operatorname{Circ}$ est $\simeq\mathbb{Z}/4\mathbb{Z}$, celui d'un élément de $\mathcal{P}_2(I)$ est $\simeq\mathbb{F}_2\times\mathbb{F}_2$) -- c'est le premier des deux, \inserted{\unsure{i.e.} $\operatorname{Circ}(I)$}, qui peut s'interpréter comme \deleted{l'application}\inserted{l'ens des} \deleted{\unsure{de $\varphi$}}\inserted{des} faces d'un cube\ldots].

L'application de degré $2$ $(*)$ définit un cube \inserted{comb.} dont l'ens. des faces est $\operatorname{Circ}(I)$, \inserted{avec l'antipodisme $u\mapsto u^{-1}$}, l'ens des \deleted{\unsure{sommets}}\inserted{arêtes} (resp. des sommets) \deleted{\unsure{est}} \inserted{est} l'ens des $12$ sections partielles \deleted{\unsure{tous}} sur des parties de card. $2$ de $J$ (resp. des \deleted{\unsure{six}} $8$ sections de $\operatorname{Circ}(I)$ sur $J$), l'incidence entre \deleted{\unsure{les faces}}\inserted{éléments de dim. $2,1,0$} étant le prolongement des sections partielles. \deleted{À \unsure{tout}} Il y aurait
\margin{Bien sûr, $G$ opère sur le cube par isomorphismes.}
% archive foliation: 22
